% ==============================================================================
% =================================== Aufgabe 3 =================================
% ==============================================================================

\chapter{Aufgabe}
\label{chap:3 Aufgabe}
\thispagestyle{fancy}
\section{Wissenschaftliche Methoden} \par Zur Überprüfung der aufgestellten Hypothese wird ein methodisches Vorgehen gewählt, das sowohl statistische Breitenwirkung als auch inhaltliche Tiefe kombiniert.
\begin{enumerate}[label=\thechapter.\arabic*), leftmargin=*]
    \item \textbf{Darstellung der wissenschaftlichen Methoden} \par Es kommen die folgenden drei Methoden zum Einsatz:
    \begin{itemize}
        \item \textbf{Quantitative Befragung (\textit{Online-Survey}):} \par Einsatz standardisierter Fragebögen zur statistischen Ermittlung von Korrelationen zwischen den Variablen „\textbf{\textit{Führungsrolle}}“ (\textit{A}) und „\textbf{\textit{Team-Motivation}}“ (\textit{B}).
        \item \textbf{Qualitative Experteninterviews:} \par Durchführung leitfadengestützter Gespräche mit Projektbeteiligten zur Exploration individueller Motivationsfaktoren im Transformationsprozess.
        \item \textbf{Systematisches Review (\textit{Literaturauswertung}):} \par Systematische Auswertung und Aggregation bereits vorhandener empirischer Studienergebnisse zur Validierung des aktuellen Forschungsstandes.
    \end{itemize}
    \item \textbf{Konkreter methodischer Weg zur Hypothesenprüfung}
        \begin{itemize}
            \item Im Rahmen der \textbf{quantitativen Erhebung} werden validierte Skalen zur Messung von Servant Leadership und intrinsischer Motivation verwendet. Die Auswertung der Daten erfolgt durch die Berechnung des Spearman-Korrelationskoeffizienten (\textit{\acs{Rho}}), um den Zusammenhang statistisch nachzuweisen.
            \item Die \textbf{qualitativen Interviews} werden mittels der qualitativen Inhaltsanalyse nach Mayring ausgewertet. Durch systematisches Codieren und Kategorisieren werden Motivationsmuster im Textmaterial identifiziert.
            \item Das \textbf{systematische Review} stützt sich auf eine gezielte Datenbankrecherche unter Verwendung definierter Inklusions- und Exklusionskriterien (\textit{\acs{z. B.} Erscheinungszeitraum der letzten 10 Jahre}), um eine verallgemeinerbare Gesamtaussage zu treffen.
        \end{itemize}
    \cleardoublepage
    \item \textbf{Tabellarische Darstellung der Quellen für die Methodik}
        \begin{table}[H]
            \centering
            \hspace*{-1.5cm}
            \begin{tabular}{|>{\raggedright\arraybackslash}p{2.6cm}|>{\centering\arraybackslash}p{0.8cm}|>{\raggedright\arraybackslash}p{3.7cm}|>{\raggedright\arraybackslash}p{4.7cm}|>{\raggedright\arraybackslash}p{4cm}|}
                \hline
                \textbf{Autor(en)} & \textbf{Jahr} & \textbf{Titel / Thema} & \textbf{Methodischer Fokus} & \textbf{Quelle / Verweis} \\
                \hline
                Preußig, J. & 2015 & Agiles Projektmanagement & Theoretische Fundierung agiler Rollenmodelle und Servant Leadership & Haufe Verlag \cite{2} \\
                \hline
                Timinger, H. & 2017 & Modernes Projektmanagement & Strukturvergleich klassischer und agiler Prozessmodelle & Wiley Verlag \cite{3} \\
                \hline
                Busse, R. et al. & 2020 & Rethinking the roles of PMM & Validierung von Reifegradmodellen mittels Kausalanalyse & Eur. J. Int. Management \cite{4} \\
                \hline
                Busse \& Regenberg & 2019 & Impact of Leadership Inclusiveness & Statistische Analyse der Korrelation zwischen Führung und Engagement & J. of Leadership \& Org. Studies \cite{5} \\
                \hline
                Busse \& Weidner & 2020 & Qualitative Investigation on EE & Qualitative Exploration von Führung in agilen Kontexten & Leadership \& Org. Dev. Journal \cite{6} \\
                \hline
                Mayring, P. & 2015 & Qualitative Inhaltsanalyse & Regelgeleitete Auswertungstechnik für Experteninterviews & Beltz Verlag \cite{7} \\
                \hline
                Saunders, M. et al. & 2016 & Research Methods for Business & Systematik des Forschungsdesigns (Research Onion) & Pearson Education \cite{8} \\
                \hline
                Busse \& Warner & 2017 & Legacy of Hawthorne experiments & Methodik der kritischen Literaturauswertung (Review-Verfahren) & History of Economic Ideas \cite{9} \\
                \hline
                Bohinc, T. & 2011 & Soft Skills für Projektleiter & Integration verhaltensbezogener "Soft Facts" in die Führung & GABAL Verlag \cite{10} \\
                \hline
                Litke, H.-D. & 2007 & Projektmanagement & Techniken und Verhaltensweisen der Projektsteuerung & Hanser Verlag \cite{11} \\
                \hline
            \end{tabular}
            \caption{Tabellarische Darstellung der methodischen Fachliteratur}
            \label{tab:literatur}
        \end{table}
    \cleardoublepage
    \item \textbf{Tabellarische Aufbereitung der Ergebnisse des systematischen Reviews}
        \begin{table}[H]
            \centering
            \hspace*{-1.5cm}
            \begin{tabular}{|>{\centering\arraybackslash}p{0.8cm}|>{\raggedright\arraybackslash}p{2.0cm}|>{\raggedright\arraybackslash}p{2.0cm}|>{\raggedright\arraybackslash}p{2.2cm}|>{\raggedright\arraybackslash}p{2.2cm}|>{\raggedright\arraybackslash}p{2.8cm}|>{\raggedright\arraybackslash}p{2.8cm}|}
                \hline
                \textbf{Jahr} & \textbf{Autor(en)} & \textbf{Zeitschrift} & \textbf{Stichprobe} & \textbf{Methode} & \textbf{Ergebnis} & \textbf{Besonderheit} \\
                \hline
                2020 & Busse, R. et al. & Eur. J. Int. Management & N=66 Projektmanager & Struktur\-gleichungs\-modell & PMM korreliert positiv mit Projekt- und Unternehmensperformance & Fokus auf deutsche Organisationen \\
                \hline
                2019 & Busse \& Regenberg & J. of Leadership \& Org. Studies & Quantitativ & Schriftliche Befragung & Inklusive Führung steigert das Mitarbeiterengagement signifikant & Neuer Erklärungsansatz für Führungsstile \\
                \hline
                2020 & Busse \& Weidner & Leadership \& Org. Dev. Journal & 10 Führungskräfte & Qualitative Interviews & Distant Leadership erfordert digitale Kollaboration für Engagement & Erprobung im agilen Umfeld \\
                \hline
                2021 & Spiegler, S. V. et al. & Empirical Software Engineering & Agile Teams & Empirische Untersuchung & Rollenverteilung zwischen Scrum Master und Team wandelt sich mit der Reife & Längsschnitt\-betrachtung \\
                \hline
                2026 & Herzing, I. M. & Saudi J Bus Manag Stud & Sekundär\-quellen & Systema\-tisches Review & Die Rolle des Scrum Masters entwickelt sich hin zu strategischen Aufgaben & Aktueller Forschungsstand 2026 \\
                \hline
            \end{tabular}
            \caption{Übersicht der empirischen Studienergebnisse zum agilen Projektmanagement}
            \label{tab:studien}
        \end{table}
    \cleardoublepage
    \textbf{\textit{Hinweis zur Quellenauswahl:}} \par \textit{Während \hyperref[tab:literatur]{Tabelle \ref{tab:literatur}} die methodische Fachliteratur (\textit{10 Quellen}) abbildet, umfasst \hyperref[tab:studien]{Tabelle \ref{tab:studien}} ausschließlich die für den systematischen Review relevanten empirischen Studien. Da Lehrbücher und theoretische Abhandlungen keine Primärdaten/Stichproben enthalten, können sie gemäß WAM04 nicht in diese Ergebnismatrix aufgenommen werden.}
\end{enumerate}

% ==============================================================================
% =================================== Aufgabe 3 =================================
% ==============================================================================